% Thorne & Campolattaro 1967, ApJ 149, 591
% Page 19 (p. 609) transcription
%
% NOTES ON EQ. (C7):
%   This equation spans six display lines. The top of the page carries a journal
%   running-head strip that partially obscures line 1. Terms marked "(?)" below
%   are uncertain due to image resolution or occlusion.

and $8(G^1 + 2G\gamma) = 8\pi(T_r^0 - T_\phi^0 + 2T_r^r)$, when combined with equation (8c), is
%
\begin{align}
H_{0,tt}
  &= -e^{-2\Lambda} H_0''
     + e^{-2\Lambda} \bigl[1 - 5m/r + 2\pi r^2 - 3p\bigr] H_0'   % "3p" possibly "3\rho"
\notag \\
  &\quad
     + 2\pi^{-1} r^{-1} \bigl\{4m/r - 10m^2/r^2 - 3p\bigr\} p    % prefactor uncertain
\notag \\
  &\quad
     + 2\pi r^{-1} \bigl[1 - \gamma + (\gamma - 3)\,2m/r\bigr]
     - 32\pi r^2 e^{-2\Lambda} p^3                                  % exponent on p uncertain
\notag \\
  &\quad
     + l[l+1]\,e^{-2\Lambda} H_0
     + 2r^{-2} \bigl[-2 + 6m/r + 8\pi r^2 \rho\bigr] K'
\notag \\
  &\quad
     + 8\pi e^2 (\rho + \gamma p) K
     - 8\pi e^2 (\rho + \gamma p)\,r^{e^{-\Lambda}} W''             % exponent on r uncertain
\notag \\
  &\quad
     - 8\pi (r^2 - p)^2\,r^{-1} e^{-\Lambda} W
     - 8\pi l(l+1)\,e^2 (\rho + \gamma p)\,r^{-1} V                % (r^2-p)^2 uncertain
\notag \\
  &= 0 \,.
\tag{C7}
\end{align}


\section*{APPENDIX D}
\subsection*{BOUNDARY CONDITIONS FOR EVEN-PARITY EIGENFUNCTIONS}

Near the center of the star ($r \ll M$, $r \ll r_s^{-1/3}$), the eigenfunctions (14) reduce to the simpler
form
%
\begin{align}
H    &= \tfrac{1}{2} r^l K'' + \bigl[1 - l(l+1)/2\bigr] K \,, \tag{D1a} \\[4pt]
H' + \bigl[1 + l(l+1)/2\bigr] H/r
     &= K'' + 3K' + \bigl[1 - l(l+1)/2\bigr] K/r \,, \tag{D1b} \\[4pt]
W'   &= -l(l+1)\,V \,, \tag{D1c} \\[4pt]
K' - H' &= \omega^2 e^{-\nu} \bigl(V + W/r\bigr) \,, \tag{D1d}
\end{align}
%
where $r_s$ is the value of $r$ at the center of the star. Equation (D1a) arises from the leading terms
in (14c) and (14b). Equation (D1d), the remaining content of (14c) and (14d), is the sum of
(14c) and (14d) used to eliminate $W$ and $V$.

Equations (D1) are a 4th-order system of linear differential equations, so there are four
linearly independent solutions. Only two of the independent solutions---those of equations
(15a)---are physically acceptable. Two unacceptable solutions correspond to the two arbitrary
constants $C$ and $D$ in
%
\begin{align}
K &= C\,r^{-(l+1)} + \cdots \,, & H &= C\,r^{-(l+1)} + \cdots \,,
\notag \\
W &= D\,r^{-(l+1)} + \cdots \,, & V &= \tfrac{1}{2} D/(l+1)\cdot r^{-(l+1)} + \cdots \,;
\tag{D2}
\end{align}
%
and the third unacceptable solution corresponds to the arbitrary constant, $E$, in
%
\begin{align}
K &= E\,r^{2-l(l+1)/4} + \cdots \,, \notag \\
H &= \tfrac{1}{2}\bigl\{1 - l(l+1)/2 - l(l+1)/2\bigr\} E\,r^{2-l(l+1)/4} + \cdots \,, \notag \\
W &= \omega^2 e^{-\nu}\{1+l\}\bigl[1 + l(l+1)/2\bigr] E\,r^{l(l+1)(l+1)/2} + \cdots \,, \notag \\
V &= -\omega^2 e^{-\nu}\bigl\{-1 - l(l+1)/2\bigr\} E\,r^{l(l+1)/2} + \cdots \,.
\tag{D3}
\end{align}

The solutions (D2) and (D3) are unacceptable for $l \geq 1$ because they lead to perturbations
in the density, the pressure, and the geometry of spacetime which diverge at $r = 0$. (Cf.\ eqs.\
[C3], [C4], and [7b].)

At the surface of the star we demand that the Lagrangian change in the pressure (eq.\ [160])
vanish and that the fluid displacement and the perturbation in the spatial geometry not diverge.
